% TikZ 图：匀强电场中的等效重力场
% Day 15：电荷、电场与电场强度
\begin{tikzpicture}[scale=0.86,>=Latex,line cap=round,line join=round,
  every node/.style={font=\small}]
  % Left: force composition on a positive charge.
  \begin{scope}[shift={(-3.1,0)}]
    \draw[blue!55,->] (-1.45,1.25) -- (1.45,1.25) node[right] {$\vect{E}$};
    \foreach \y in {0.62,-0.62} {
      \draw[blue!55,->] (-1.45,\y) -- (1.45,\y);
    }
    \filldraw[red] (0,0) circle (3pt);
    \node[red,above left=4pt] at (0,0) {$+q$};
    \draw[very thick,red!80!black,->] (0,0) -- (1.15,0);
    \node[red!80!black,above=4pt] at (0.78,0) {$qE$};
    \draw[very thick,black,->] (0,0) -- (0,-1.15) node[left] {$mg$};
    \draw[very thick,orange!85!black,->] (0,0) -- (1.15,-1.15)
      node[below right] {$mg'$};
    \draw[dashed,gray] (1.15,0) -- (1.15,-1.15) -- (0,-1.15);
  \end{scope}

  % Right: pendulum aligns with the resultant force.
  \begin{scope}[shift={(2.8,0)}]
    \coordinate (O) at (0,1.35);
    \coordinate (P) at (0.82,-0.65);
    \filldraw[black] (O) circle (1.4pt);
    \draw[dashed,gray] (O) -- (0,-1.15);
    \draw[thick] (O) -- (P);
    \filldraw[red] (P) circle (4pt);
    \node[red,above left=4pt] at (P) {$+q$};
    \draw[blue!55,->] (-1.45,0.88) -- (1.45,0.88) node[right] {$\vect{E}$};
    \draw[blue!55,->] (-1.45,0.25) -- (1.45,0.25);
    \draw[very thick,orange!85!black,->] (P) -- ++(0.60,-1.10)
      node[right] {$mg'$};
    \draw[very thick,black,->] (P) -- ++(0,-0.95) node[left] {$mg$};
    \draw[very thick,red!80!black,->] (P) -- ++(0.90,0);
    \node[red!80!black,above=4pt] at ($(P)+(0.66,0)$) {$qE$};
    \draw[->] (0,0.66) arc[start angle=-90,end angle=-68,radius=0.58];
    \node at (0.30,0.42) {$\theta$};
  \end{scope}
\end{tikzpicture}
